\documentclass[11pt]{article}
\usepackage[margin=2.5cm]{geometry}
\usepackage{amsmath,amssymb,booktabs,graphicx,xcolor,hyperref}

% Inline reminders. \todo = something still to be done or verified; \claim = a statement that may be made ONLY after the condition in the note is met.
\newcommand{\todo}[1]{{\color{red}[TODO: #1]}}
\newcommand{\claim}[1]{{\color{blue}[CLAIM only if: #1]}}
\newcommand{\MQ}{\mathcal{M}_Q}

\title{Validating black-hole particles in Schr\"odinger--Poisson simulations:\\
an exact finite-box benchmark and a momentum-conserving force on the particle}
\author{M.~\v{S}m\'id \and \todo{co-authors / affiliation}}
\date{\today\ --- working draft, numbers from Koios P100 runs of 2026-10-29/30}

\begin{document}
\maketitle

\begin{abstract}
\todo{Write last, 150 words.} Skeleton: (i) we derive the exact linear-theory drag on a Gaussian-softened massive particle moving through a uniform ultralight-dark-matter (ULDM)
background in a finite periodic box with switch-on at $t=0$ and $\psi$ self-gravity (discrete mode sum); (ii) we show that a point-receiver force on the particle does not conserve the total momentum
of $\psi$ plus particle, with an error growing as $(\epsilon/\lambda)^2$, and that smoothing the force with the same Gaussian as the particle's source restores it (error $\lesssim 1\,\%$, falling as $N^{-2}$);
(iii) we compare three independent split-step codes at $\epsilon=0.29\lambda$ and give a validity rule for the softening and resolution; (iv) we reproduce the stopping time of a black hole released in a soliton
(arXiv:2403.09038) to $\sim 2\,\%$ and quantify its sensitivity to the numerical parameters (\todo{only after T6}).
\end{abstract}

\section{Introduction}
\begin{itemize}
\item ULDM/FDM, Schr\"odinger--Poisson (SP) with a Fourier split-step scheme; black holes (BH) as gravitating particles or sinks; dynamical friction in ULDM is oscillatory and not captured by the classical formula
(Hui et al.\ 2017; Lancaster et al.\ 2020, arXiv:1909.06381; Boey et al.\ 2024, arXiv:2403.09038). \todo{add Wang \& Easther 2022, Buehler \& Desjacques (2207.13740), Boey et al.\ 2025 (2504.16348) after reading them; second-hand so far}
\item The numerical treatment of the particle differs between codes (point receiver with trilinear or spectral interpolation, Plummer or Gaussian softening, periodic or non-periodic particle potential); results are rarely tested against an exact reference.
\item Contributions (candidate list, see Sec.~\ref{sec:claims}).
\claim{a literature search shows that no earlier paper gives an exact finite-box benchmark and a momentum-balance analysis of the receiver}
\end{itemize}

\section{Method}
\subsection{Equations and scheme}
SP equations in kpc, $M_\odot$, Gyr; $\tilde\hbar=\hbar/m$ ($7.843\,\mathrm{kpc^2/Gyr}$ for $m=2.5\times10^{-22}$\,eV, $1.961$ for $10^{-21}$\,eV); periodic box, spectral Poisson solver with the mean density removed;
second-order kick--drift--kick split step, $\psi$ in double precision (complex64 gave a norm drift proportional to the number of steps: $-0.3\,\%$/Gyr at $h=1.25\times10^{-4}$, $N=256$).
Stability limit used by the code: $h_{\max}=\frac12\cdot\frac{4}{3\pi}\,\Delta x^2/\tilde\hbar$ (now an optional hard error, \texttt{sim.strict\_time\_step}).
\subsection{Black-hole particle (sink)}
Gaussian source of width $\epsilon$ (density $\rho_G=M(2\pi\epsilon^2)^{-3/2}e^{-r^2/2\epsilon^2}$), analytic potential, CIC deposit convolved with $e^{-k^2\epsilon^2/2}$.
Two receivers: \emph{point} (force at the particle position, trilinear interpolation) and \emph{smooth} (force filtered with the same Gaussian before interpolation; action $=$ reaction on the grid; unit test error $2\times10^{-15}$).
\subsection{Quantum Mach number}
$\MQ=\tilde\hbar v/(GM)=\lambda/(GM/v^2)$, $\lambda=\tilde\hbar/v$. Linear regime $\MQ\gg1$; $\MQ\sim1$ strongly quantum; $\MQ\ll1$ nearly classical.

\section{Exact linear benchmark for a finite periodic box}
\label{sec:theory}
\todo{Derive in an appendix (the derivation is in the docstring of tests/sink\_friction.py).} Mode equation
$\ddot u_k+\Omega_k^2u_k=-\psi_0k^2\Phi_k\,e^{-i\mathbf{k}\cdot\mathbf{v}t}$, $\Omega_k^2=\tilde\hbar^2k^4/4-4\pi G\rho_0$, Gaussian $\Phi_k$ of width $\epsilon$,
force obtained by the sum over the discrete modes of the box. Drag coefficient $C=F/F_{\rm rel}$, $F_{\rm rel}=4\pi\rho_0(GM/v)^2$.
Comparison references: (a) this exact finite-box theory (valid for $\MQ\gg1$), (b) steady-state point-source formula, Lancaster et al.\ Eq.~(4.21), tending to Eq.~(4.26) for large $\MQ$ (our own implementation; \todo{state clearly that it is ours}).

\section{Test set-up}
Uniform $\psi$, $\rho_0 M\le 3\times10^{12}$ (so the particle slows by only a few percent); box $[-9.6,9.6]^3$\,kpc, $v=30$\,kpc/Gyr, half a box of travel; $m=2.5\times10^{-22}$\,eV, $\lambda=0.2615$\,kpc;
$N=128$--$320$ ($\Delta x=0.150$--$0.060$\,kpc); $\epsilon=0.29\lambda=0.075$\,kpc fixed in physical units (0.5/0.75/1.0 cell at $N=128/192/256$), plus $0.14\lambda$ and $0.07\lambda$ at $N=256,320$;
$\MQ=17.4,\ 1,\ 0.3$ ($M=3.0\times10^6,\ 5.23\times10^7,\ 1.74\times10^8\,M_\odot$). Mean $C$ taken for $vt/\lambda>7$.

\section{Results}
\subsection{Convergence and the exact benchmark}
Table of mean $C$ against $N$ and $h$ for the smooth and the point receiver ($N=128,192,256,320$), exact linear theory $3.174$--$3.169$ for $\MQ=17.4$ (box), steady point source $3.268$.
Facts to report (T1/T2): simulation = linear theory within $0.8\,\%$ for $\MQ\ge3$; time step $h$ changes $C$ by $\le0.5\,\%$ ($N=128$, $h=4\times10^{-4}/2\times10^{-4}/10^{-4}$);
$N=128\to256$ shifts $C$ by $+0.5\,\%$. \todo{figure: $C(vt/\lambda)$ for all $N$ with the exact curve.}
\subsection{Momentum and energy balance of the receiver}
Momentum gained by $\psi$ vs lost by the particle, $P_\psi/P_{\rm lost}-1$ ($N=192$; $\MQ=17.4/1/0.3$):
point $-2.5/-6.5/-9.4\,\%$ (ours), $-3.0/-7.5/-8.7\,\%$ (jaxion); smooth $+0.5/+1.0/+0.8\,\%$ (ours, falls $\sim N^{-2}$), $0.00/0.00/-0.01\,\%$ (jaxion, spectral).
Softening scan at $\MQ=17.4$ (point): $-13/-8/-1.9/+0.3\,\%$ for $\epsilon/\lambda=0.86/0.57/0.29/0.14$. Energy: $\max|E-E_0|/K_{\rm lost}$ point $1.7$--$8\,\%$, smooth $0.26$--$1.9\,\%$.
\todo{explain the sign of the point error; caveat: our smooth receiver is not exactly reciprocal when $\epsilon\sim0.25$ cell ($+3.9\,\%$ at $N=256$); the CIC explanation is a hypothesis, untested.}
\subsection{Cross-code comparison}
\begin{table}[h]\centering\small
\begin{tabular}{lcccccc}\toprule
 & \multicolumn{2}{c}{$\MQ=17.4$} & \multicolumn{2}{c}{$\MQ=1$} & \multicolumn{2}{c}{$\MQ=0.3$}\\
 code, receiver & $N{=}192$ & bal.\ & $N{=}192$ & bal.\ & $N{=}192$ & bal.\\\midrule
 ours, point        & 3.185 & $-2.5$ & 2.820 & $-6.5$ & 2.172 & $-9.4$\\
 jaxion, point      & 3.202 & $-3.0$ & 2.860 & $-7.5$ & 2.187 & $-8.7$\\
 PyUL\_NBody, point & 3.185 & $+4.4$ & 2.869 & $+1.1$ & 2.292 & $-0.4$\\
 ours, smooth       & 3.102 & $+0.5$ & 2.647 & $+1.0$ & 2.007 & $+0.8$\\
 jaxion, smooth     & 3.119 & $0.00$ & 2.675 & $0.00$ & 2.021 & $-0.01$\\\midrule
 steady point source & \multicolumn{2}{c}{3.268} & \multicolumn{2}{c}{2.778} & \multicolumn{2}{c}{1.919}\\\bottomrule
\end{tabular}
\caption{Mean drag coefficient for $vt/\lambda>7$, $\epsilon=0.287\lambda$, and momentum balance in \% (end of run). PyUL balance is $p_{\rm end}/{\rm lost}-1$ as reported by the agent. \todo{N=256 and N=320 columns; check PyUL's balance definition.}}
\end{table}
Points to state honestly: (i) jaxion (0.0.13, unmodified, spectral double precision) and our code agree to $0.2$--$1.5\,\%$ (smooth) and $0.4$--$2.6\,\%$ (point; one outlier $-5\,\%$, $N=128$, $\MQ=0.3$);
(ii) PyUL\_NBody (non-periodic particle potential, zero force within ~4 cells of the edge) agrees to $0.3\,\%$ at $\MQ=17.4$, $0.7$--$2\,\%$ at $\MQ=1$ and $+5.5$--$5.9\,\%$ at $\MQ=0.3$, $N\ge192$ (cause not found);
(iii) the comparison is not blind (the reference numbers were given to the agent running jaxion), the three codes share the periodic spectral split-step family, and the particle kernels differ ($e^{-k^2\epsilon^2/2}$ without CIC in jaxion, Gaussian $\times$ CIC in ours, Plummer in PyUL).
\subsection{The regime $\MQ\lesssim1$}
$C(t)$ oscillates; the oscillation phase agrees between jaxion and our code (correlation $0.96$--$1.00$ for $N\ge192$ at $\epsilon=0.29\lambda$, $0.90$ for smooth at $0.07\lambda$).
The mean exceeds the steady-state value by $5$--$13\,\%$ at $\MQ=0.3$ at every $N$ and travel length tested; with $\epsilon\to0$ ($N=256$, smooth) the mean stays near $2.10$ (jaxion), $2.05$--$2.09$ (ours), against the steady state $1.919$.
\claim{an independent reference (second code or a published nonlinear result) agrees at $N=320$ and for a longer travel; until then this regime is ``converged in $N$ and $h$, no exact reference''.}
\subsection{Validity rule}
Point receiver: $\epsilon\gtrsim0.5$ cell and $\epsilon/\lambda\lesssim0.3$, $\lambda\ge3$ cells (below $\sim0.5$ cell the phase of $C(t)$ and the momentum balance break down: correlation $0.02$--$0.24$, balance up to $+230\,\%$ in jaxion).
Smooth receiver: usable to $\sim0.25$ cell with a few-percent momentum error. \todo{derive the rule from the full $(\MQ,\epsilon/\lambda,N)$ map, not from these points.}

\section{A black hole in a soliton (arXiv:2403.09038 set-up)}
Set-up: soliton $5.57\times10^7M_\odot$ (free ground state), $m=10^{-21}$\,eV, BH $4\times10^6M_\odot$ at rest at $(0.125,0.125,0)$\,kpc, box $2$\,kpc, $N=192$, $\epsilon=0.5$ cell, $h=1.06\times10^{-5}$\,Gyr ($0.9\,h_{\max}$).
Result (one run, 60\,Myr): peak speed $79.6$\,km/s (paper $\sim80$); stopping time (minimum of the running-maximum envelope of the speed) $45.9$\,Myr against $A/M+B=46.9$\,Myr ($A=1.38\times10^8\,M_\odot$\,Myr, $B=12.4$\,Myr); BH re-accelerates after $\sim50$\,Myr.
Numerical diagnostics of the same run: $E_{\rm total}$ drift $0.6$--$0.8\,\%$, total momentum of $\psi$+BH grows to $\sim6\times10^6\,M_\odot$\,kpc/Gyr ($2$--$6\,\%$ of the exchanged momentum).
\todo{T6: dependence of $t_{\rm stop}$ and of the onset of reheating on $\epsilon/\Delta x\in\{0.25,0.5,1,1.5\}$, $N\in\{128,192,256\}$, point vs smooth, box size, sponge vs periodic, frozen soliton; 150\,Myr run for the reheating.}
\todo{an independent code on the same set-up: the PyUL\_NBody attempt gave $t_{\rm stop}=8.2$\,Myr (standard PyUL profile, unrelaxed, periodic box, own $t_{\rm stop}$ definition) and is not a conclusion.}
\claim{T6 shows the reheating onset is (not) sensitive to numerical parameters; the authors of 2403.09038 themselves write that at late times results become dominated by numerical effects.}

\section{Discussion and limitations}
Periodic box (mean-density term $-(4\pi/3)Gmr/V$, usable to $r\lesssim L/4$); uniform background only for the benchmark; Gaussian/point BH, no accretion or gas in the validated part; $\epsilon\to0$ at $\MQ\le1$ unconverged;
no GPU/AMR scaling claims (P100 at $N=320$ about 14--18\,min per friction case; $N=384$ exceeds 16\,GB); sponge boundaries (complex64) not tested.

\section{Conclusions}
\todo{After the T6 and $N=320$ results.}

\section{Reproducibility}
Code and tests (\texttt{tests/run\_all.py}, regression baselines bitwise identical on a P100, friction scans, T5 scripts), scripts for all figures, commit hash \todo{fill}, DOI of an archive \todo{Zenodo}.
\todo{licence, CITATION.cff, README (installation, minimal example, units, restart) before submission.}

\section{What may and may not be claimed}\label{sec:claims}
\begin{itemize}
\item \textbf{May (supported now):} exact finite-box linear benchmark reproduced to $<1\,\%$ for $\MQ\ge3$; point receiver violates momentum conservation with a size depending on $\epsilon/\lambda$ and $\MQ$ in two codes; smooth receiver restores it;
two codes agree to $\sim1\,\%$ at $\epsilon=0.29\lambda$; our code reproduces the $t_{\rm stop}$ of 2403.09038 to $\sim2\,\%$ in one run.
\item \textbf{May not (yet):} novelty (literature not read in full); a validation for $\MQ\lesssim1$ (no exact reference; PyUL disagrees by $6\,\%$ at $\MQ=0.3$); that the reheating is physical or numerical; that the codes are independent in the strong sense.
\item \textbf{Venues:} RNAAS (short note), Astronomy \& Computing, JOSS (as a companion software paper).
\end{itemize}

\appendix
\section{Derivation of the finite-box mode sum} \todo{from tests/sink\_friction.py}
\section{Parameter tables for all runs} \todo{from tests/t2\_aggregate.py output, jaxion table.md, PyUL analysis\_table.md}

\begin{thebibliography}{9}
\bibitem{Lancaster2020} L.~Lancaster, C.~Giovanetti, P.~Mocz, Y.~Kahn, M.~Lisanti, D.~N.~Spergel, JCAP 01 (2020) 001, arXiv:1909.06381.
\bibitem{Boey2024} arXiv:2403.09038 \todo{full author list and journal reference from the arXiv listing; I only verified the content of the PDF, not the metadata}.
\bibitem{Hui2017} L.~Hui, J.~P.~Ostriker, S.~Tremaine, E.~Witten, Phys.\ Rev.\ D 95, 043541 (2017).
\bibitem{Edwards2018} F.~Edwards, E.~Kendall, S.~Hotchkiss, R.~Easther, JCAP 10 (2018) 027 (PyUltraLight).
\bibitem{Schive2014} H.-Y.~Schive, T.~Chiueh, T.~Broadhurst, Nature Physics 10, 496 (2014).
\bibitem{Wang2022} Y.~Wang, R.~Easther, Phys.\ Rev.\ D 105, 063523 (2022).
\bibitem{jaxion} \todo{reference of jaxion (repository JaxionProject/jaxion, version 0.0.13, commit ef96a94b) -- find the paper}
\bibitem{PyULNBody} \todo{reference of PyUL\_NBody (Sifyrena, commit 879b8e7e)}
\end{thebibliography}
\end{document}
